% Cálculo Vetorial: 2 vectores 
% Faro, sáb 17 jan 2026 20:49:39 WET
% Written by Tordar 
% pstricks2pdf.sh 2vectores
\documentclass{article}

\pagestyle{empty}

\usepackage{graphicx}
\usepackage{pst-all} 
\usepackage{pstricks-add}
\usepackage{pst-slpe}

\input{apnum}

\begin{document}

% Let's go:
\begin{pspicture}(0,0)(10,10)
% Grelha: 
    \psgrid[xunit=1, yunit=1, 
            subgriddiv=1, 
            gridlabels=10pt, % font size for numbers
            gridwidth=0.5pt,
            griddots=10,
            subgridwidth=0.2pt,
            subgridcolor=lightgray,
            gridcolor=gray](-3,-1)(3,4)
\SpecialCoor
% Vectores: 
\psline[linewidth=3pt,linecolor=blue]{->}(2.5;30)
\psline[linewidth=3pt,linecolor=blue]{->}(3.8;120)
\psline[linewidth=3pt,linecolor=red ]{<->}(2.5;30)(3.8;120)
\pcline[offset=-12pt,linewidth=2pt]{<->}(2.3;30)(3.8;120)\ncput*{\scalebox{1.2}{$r$}}
% Ângulos:
\psarc[linewidth=0.8mm,linecolor=black,linestyle=dashed,dash=4mm 1mm]{->}(0,0){1.5}{0}{30}
\psarc[linewidth=0.8mm,linecolor=black,linestyle=dashed,dash=4mm 1mm]{->}(0,0){1}{0}{120}
% Texto: 
\rput{0}(2;15){\scalebox{1}{$30^{\circ}$}}
\rput{0}(1.5;60){\scalebox{1}{$120^{\circ}$}}
\end{pspicture}

\end{document}
